\section{Wide separations at a fixed original scale}\label{sec:rounds}

The local estimate pays for a metric change by a decrease in a trace.
That decrease can be summed over arbitrarily many labels. We first carry
out this summation, then explain why different labels may require
different finite numbers of rounds.

Fix \(L_*=1.004\), and choose \(\delta>0\), depending only on \(n\), so small that
\begin{equation}\label{rounds:delta}
  2(n+1)\delta<1-L_*^{-2}.
\end{equation}
Let \(W,T\) be the constants in \Cref{step:main}. Enlarge \(T\) if necessary
to make \(T\ge1\), and choose
\begin{equation}\label{rounds:original-scale}
  R_0\ge\sqrt2\,T,\qquad D=WR_0.
\end{equation}
All of these constants depend only on \(n\).

\begin{lemma}[A complete round]\label{rounds:round}
Suppose that \(h_0\) is a complete metric, \(t_0\ge T\) is smooth, and
\[
 \mathcal D(h_0,t_0)\ge c+t_0^2
\]
for a constant \(c\). Let \((A_i,B_i)\) be a finite or countable family of
nonempty closed pairs, with each \(A_i\) compact and
\(\dist_{h_0}(A_i,B_i)\ge W\). Suppose there are locally finite compact
sets \(K_i\) such that the open \(h_0\)-distance-\(W\) neighborhood of
\(A_i\) is contained in \(K_i\setminus B_i\).
Then there are smooth \(h_1\ge h_0\) and \(t_1\ge t_0\), obtained by a
locally finite sequence of compactly supported changes inside the
\(K_i\), such that
\begin{align}
 \dist_{h_1}(A_i,B_i)&\ge L_*W \quad\text{for every }i,\label{rounds:separation}\\
 \mathcal D(h_1,t_1)&\ge c+
 \bigl(1-(n+1)\delta\bigr)t_1^2.\label{rounds:curvature}
\end{align}
The metric \(h_1\) is complete.
\end{lemma}

\begin{proof}
Enumerate the labels and choose \(d_i>0\) with
\(\sum_i d_i<\delta\). Apply \Cref{step:main} in this order, keeping
\(h_0\) as the background metric throughout. Earlier changes increase
distances, so every pair still meets the separation hypothesis when it
is processed. If \((h^{(j)},t^{(j)})\) denotes the data after \(j\) changes,
write
\[
 Q_j=\tr_{h^{(j)}}h_0,\qquad
 A_j=1-\delta(n-Q_j)-\sum_{i\le j}d_i.
\]
Thus \(Q_0=n\), \(0<Q_j\le n\), and \(A_j\le1\).
We claim that
\begin{equation}\label{rounds:trace-induction}
 \mathcal D(h^{(j)},t^{(j)})\ge c+A_j(t^{(j)})^2.
\end{equation}
It holds at \(j=0\). In a single step, abbreviate the old coefficient
by \(A\), the new one by \(A'\), and the old and new functions by
\(t,t'\). The local estimate gives
\begin{align*}
 \mathcal D(h',t')
 &\ge c+At^2+(t'^2-t^2)
       -\bigl(\delta(Q-Q')+d\bigr)t'^2\\
 &=c+A't'^2+(1-A)(t'^2-t^2)\ge c+A't'^2.
\end{align*}
The last inequality uses both \(A\le1\) and \(t'\ge t\ge T>0\).
This proves \eqref{rounds:trace-induction}. Its coefficient is greater
than \(1-(n+1)\delta\), which proves \eqref{rounds:curvature}.
Each processed separation satisfies \eqref{rounds:separation}, and
later changes preserve it.

The supports lie in the locally finite family \((K_i)\). Hence every
point has a neighborhood on which the sequence is eventually constant;
its limit is smooth and satisfies the same inequalities. Completeness
follows from \(h_1\ge h_0\): an \(h_1\)-Cauchy sequence converges in
\(h_0\), and near its limit the two smooth metrics are equivalent.
\end{proof}

The passage from a positive scalar-curvature form to circle warps uses
torical symmetrization; see \citet[Section~2]{gromov2018Metric}. The
following identity fixes the finite-dimensional convention used here.

\begin{lemma}[Scalar curvature of diagonal circle warps]\label{rounds:warp-identity}
For smooth functions \(p_1,\ldots,p_k\) on \((M,h)\), the metric
\(\widehat h=h+\sum_{\alpha=1}^k e^{2p_\alpha}\dd\theta_\alpha^2\)
on \(M\times\T^k\) has scalar curvature
\[
 \Scal_{\widehat h}=\Scal_h-2\Delta_h P-|\dd P|_h^2
                    -\sum_{\alpha=1}^k|\dd p_\alpha|_h^2,
 \qquad P=\sum_{\alpha=1}^k p_\alpha.
\]
\end{lemma}

\begin{proof}
Let \(E_\alpha=e^{-p_\alpha}\partial_{\theta_\alpha}\) be the unit
vertical fields. For horizontal fields \(X,Y\), the Koszul formula gives
\[
 \widehat\nabla_XY=\nabla^h_XY,\quad
 \widehat\nabla_XE_\alpha=0,\quad
 \widehat\nabla_{E_\alpha}X=(Xp_\alpha)E_\alpha,\quad
 \widehat\nabla_{E_\alpha}E_\alpha=-\grad_h p_\alpha,
\]
and \(\widehat\nabla_{E_\alpha}E_\beta=0\) for \(\alpha\ne\beta\).
Therefore, for a horizontal unit vector \(X\), the mixed sectional
curvatures are
\[
 K(X,E_\alpha)=-\Hess_h p_\alpha(X,X)-(Xp_\alpha)^2,
 \qquad
 K(E_\alpha,E_\beta)=-\langle\dd p_\alpha,\dd p_\beta\rangle_h.
\]
Summing twice over all orthonormal two-planes gives the identity.
\end{proof}

\begin{theorem}[Prescribed widths for a locally finite family]\label{rounds:wide}
Let \((M,g)\) be complete, connected and smooth, with
\(\Scal_g\ge1\). There is a countable set of centers \((x_i)\) such that
\[
 A_i=\overline B_g(x_i,D),\qquad
 B_i=M\setminus B_g(x_i,2D)
\]
have the following properties: the \(A_i\) cover \(M\), and the closed
balls \(\overline B_g(x_i,2D)\) are locally finite. Suppose that every
\(B_i\) is nonempty. Given arbitrary finite numbers \(L_i>0\), there
exist a complete smooth metric \(h\ge g\) and a countable family of
smooth functions \((p_\alpha)\), whose supports form a locally finite
family of compact sets, for which
\begin{align}
 \dist_h(A_i,B_i)&>L_i \quad(i\ge1),\label{rounds:prescribed}\\
 \Scal_h-2\Delta_h\sum_\alpha p_\alpha
       -\left|\dd\sum_\alpha p_\alpha\right|_h^2
       -\sum_\alpha|\dd p_\alpha|_h^2&\ge\frac25.\label{rounds:stabilized}
\end{align}
Every sum in \eqref{rounds:stabilized} is locally finite. If some
\(B_i\) is empty, then \(\diam_g M\le4D\).
\end{theorem}

\begin{proof}
A complete connected Riemannian manifold is proper. Choose a maximal
\(D\)-separated set of centers. Maximality gives the covering by closed
\(D\)-balls. A compact set contains only finitely many centers; more
generally, centers whose closed \(2D\)-balls meet a fixed compact set
lie in a fixed compact neighborhood of that set. This proves local
finiteness of the outer balls and countability of the net. The triangle
inequality gives \(\dist_g(A_i,B_i)\ge D\). The last assertion follows
immediately if \(M=B_g(x_i,2D)\).

We keep a scalar reserve of \(1/2\) in the original units. Put
\(R_j=L_*^jR_0\), start with \(h^{[0]}=g\), \(t^{[0]}=T\), and work in
round \(j\) with the metric \(h^{[j]}/R_j^2\). At the start of that
round we require
\begin{equation}\label{rounds:start}
 \mathcal D(h^{[j]}/R_j^2,t^{[j]})
 \ge \frac{R_j^2}{2}+(t^{[j]})^2.
\end{equation}
For \(j=0\), this follows from \(\Scal_g\ge1\) and
\(R_0^2/2\ge T^2\).

A label remains active until its separation exceeds its prescribed
\(L_i\). At the start of round \(j\), every active pair has separation
at least \(WR_j\) in \(h^{[j]}\); this is initially true and follows
inductively from \eqref{rounds:separation}. Apply \Cref{rounds:round} to
the active pairs in the working metric, with reserve \(c=R_j^2/2\).
The required support condition follows from
\(\dist_{h^{[j]}/R_j^2}(A_i,B_i)\ge W\): the open distance-\(W\)
neighborhood of \(A_i\) misses \(B_i\), and is therefore contained in
the original outer ball \(B_g(x_i,2D)\). We may take
\(K_i=\overline B_g(x_i,2D)\).

Return the output metric to original units and denote it by
\(h^{[j+1]}\). It dominates \(h^{[j]}\). Scalar curvature, the
Laplacian and squared gradient norms all multiply by \(R^2\) when a
metric is divided by \(R^2\). Consequently, on passing from working
scale \(R_j\) to \(R_{j+1}\), the round estimate becomes
\[
 \mathcal D(h^{[j+1]}/R_{j+1}^2,t^{[j+1]})
 \ge \frac{R_{j+1}^2}{2}
 +L_*^2\bigl(1-(n+1)\delta\bigr)(t^{[j+1]})^2.
\]
By \eqref{rounds:delta}, the coefficient of the final square is greater
than one. Thus \eqref{rounds:start} is restored. In original units the
reserve remains exactly \(1/2\); it has not been spent by the change
of scale.

Each fixed label becomes inactive after finitely many rounds, since
its separation is at least \(WR_j\) as long as it remains active and
\(R_j\to\infty\). A compact set meets only finitely many of the
original outer balls. Each of their labels is processed only finitely
many times. Thus there are only finitely many changes on a neighborhood
of any compact set over the \emph{entire} construction, not merely in
each round. The limiting data \((h,t)\) are smooth, satisfy
\[
 h\ge g,\qquad t\ge T,\qquad \mathcal D(h,t)\ge\frac12,
\]
and satisfy \eqref{rounds:prescribed}. The same comparison argument as
in \Cref{rounds:round} proves completeness of \(h\).

It remains to realize the drift term by ordinary finite-dimensional
warped products on each compact window. Enumerate the nonzero compactly
supported increments of \(t\) as \(\tau_\nu\ge0\), so that
\(t-T=\sum_\nu\tau_\nu\). The supports are locally finite. For each
\(\nu\), choose an integer \(m_\nu\ge1\) such that
\[
 \frac1{m_\nu}\sup_M|\dd\tau_\nu|_h^2\le\frac{2^{-\nu}}{10}.
\]
The supremum is finite because the support is compact. Use \(m_\nu\)
fresh log warps, all equal to \(-\tau_\nu/m_\nu\). Listing these
functions as \((p_\alpha)\), we obtain
\[
 \sum_\alpha p_\alpha=-(t-T),\qquad
 \sum_\alpha|\dd p_\alpha|_h^2\le\frac1{10}.
\]
On a neighborhood meeting only finitely many supports, equip a finite
product of circles with metric
\(h+\sum_\alpha e^{2p_\alpha}\dd\theta_\alpha^2\), omitting all unused
unit factors. Its scalar curvature is the left side of
\eqref{rounds:stabilized}; this is the identity in \Cref{rounds:warp-identity}. Substitution gives
\(\mathcal D(h,t)-\sum_\alpha|\dd p_\alpha|_h^2\ge1/2-1/10=2/5\).
This proves \eqref{rounds:stabilized}.
\end{proof}

\begin{remark}\label{rounds:no-infinite-torus}
The list of warps is bookkeeping for compatible finite-dimensional
calculations. We never construct an infinite-dimensional torus.
Adding unit factors changes neither the weight nor the scalar formula.
After the radial partition in \Cref{countable:construction}, only
finitely many initial warp supports occur on any whole packet; later
stages add finite groups of factors. This stronger finiteness is used
in the one-dimensional obstruction, not just local finiteness of a
series on the original manifold.
\end{remark}
